\section{About the Topic}

\subsection{Key Terms}

The following terms will be mentioned various times throughout the study guide and must be known by delegates in order to better understand and address the problem of child soldier use and recruitment.

\subsubsection{Child Soldiers}

Any person under 18 years of age who is recruited, used, or associated with government armed forces of non-state armed forces groups in any capacity. Note that this definition does not only refer to children who have been directly engaged in fighting, but also includes those who have been used to serve as messengers, cooks, porters, spies, or sexual slaves\footnote{United Nations Office of the Special Representative of the Secretary-General for Children and Armed Conflict. \emph{Child Recruitment and Use}. \url{https://childrenandarmedconflict.un.org/en/six-grave-violations/child-soldiers}}.

\subsubsection{Disarmament}

The laying aside of arms or the reduction of military capacity to a set minimum\footnote{\emph{Disarmament}. Merriam-Webster. \url{https://www.merriam-webster.com/dictionary/disarmament}}.

\subsubsection{Demobilization}

The formal act and controlled discharge of active combatants from armed forces or other groups\footnote{Social Development Department. \emph{Disarmament, Demobilization and Reintegration}. No. 119. 2009. \url{https://au.int/sites/default/files/documents/39242-doc-192._social_development_notes._conlfict_crime_and_violence.pdf}}.

\subsubsection{Reintegration}

The process through which minors formerly related to armed groups or conflicts are brought back to society by reunifying them with their families and providing them with physical and psychological assistance, economic support, education, or vocational training.

\subsubsection{Paris Principles (2007)}

A set of international guidelines which seek to protect children in war contexts, prevent the re-recruitment of child soldiers, and help them reintegrate back into society.

\subsubsection{Optional Protocol on Children in Armed Conflict (OPAC)}

Adopted in the year 2000, this protocol builds upon the statements made in the famous Convention on the Rights of the Child by further protecting children from being recruited by armed groups to participate in hostile environments and by setting the minimum age for conscription at 18.

\subsubsection{Grave Violations framework}

In the year 2005, the UNSC created a Monitoring and Reporting Mechanism (MRM) which sought to monitor, document and report on violations against children’s rights taking place all around the world. It recognised six types of grave violations of children’s rights: the killing and maiming of children; the recruitment or use of children as soldiers; the commission of sexual violence against children; the execution of child abductions; the performance of attacks against schools or hospitals; and the denial of humanitarian access for children\footnote{United Nations Office of the Special Representative of the Secretary-General for Children and Armed Conflict. \emph{The Six Grave Violations}. \url{https://childrenandarmedconflict.un.org/en/six-grave-violations}}.

\subsubsection{Psychosocial Support (PSS)}

This relates to the rehabilitation efforts performed in an attempt to address both the psychological and social impacts of war on former child soldiers. The most common psychological wounds that arise in these contexts are trauma or PTSD, while the main social challenges that children formerly associated with armed conflict experience usually range from stigma, to loss of educational and professional opportunities.

\subsection{Introduction to the Topic}

The illicit recruitment and use of children in armed conflict illustrate a serious, if not one of the gravest violations of international humanitarian law.\footnote{History Extra. \emph{Innocence under Fire: How Children were Used in Warfare of the Past.} 2021.\url{https://www.historyextra.com/period/first-world-war/how-children-used-warfare-fighting-child-soldiers/}}  Through history and in present reality, tens of thousands of children – many of whom are under the age of 15 – are faced with being forcibly conscripted, abducted or coerced into service by both state militaries and non-state armed groups. They are often used as fighters, spies, or human shields, and are stripped of their childhood and dignity at a very young age. The persistence of the recruitment of child soldiers, despite decades of international condemnation, reflects the structural failures for accountability and post-conflict support.

As a potential solution to address the aftermath of child recruitment, the Disarmament, Demobilization, and Reintegration (DDR) emerged.\footnote{United Nations. \emph{Integrated Disarmament, Demobilization and Reintegration Standards (IDDRS).} 2014.\url{https://www.unddr.org/iddrs/}} The DDR attempts to break the cycle of conflict by facilitating the transition of former combatants, particularly children, back into civilian society. Such a framework addresses concerns that are not usually considered in military solutions that focus on immediate security. Each phase of the DDR attempts to deal with distinct challengers: disarmament being designed to include children who have never bore arms; demobilization requiring careful identification and interim care; and reintegration, which demands sustained long-term investment in children, families, communities and individuals that the system has failed to protect.

International frameworks already in place governing child soldiers are abundant but remain fragmented in terms of their enforcement. The Convention on the Rights of Child (1989),\footnote{United Nations General Assembly. \emph{Convention on the Rights of the Child.} UN Doc. A/RES/44/25. November 20, 1989} the Optional Protocol on Children in Armed Conflict (2000)\footnote{United Nations General Assembly. \emph{Optional Protocol to the Convention on the Rights of the Child on the Involvement of Children in Armed Conflict.} UN Doc. A/RES/54/263. May 25, 2000}, the Rome Statute (1998)\footnote{United Nations Diplomatic Conference of Plenipotentiaries. \emph{Rome Statute of the International Criminal Court.} UN Doc. A/CONF.183/9. July 17, 1998}, and the Paris Principles (2007)\footnote{UNICEF. \emph{The Paris Principles: Principles and Guidelines on Children Associated with Armed Forces or Armed Groups.} February 2007.\url{https://www.unicef.org/emerg/files/ParisPrinciples310107English.pdf}} all collectively prohibit the recruitment and the use of children, state the obligations of each state and non-state actors, and establish compliance procedures for DDR programming. The United Nations Security Council has also passed many resolutions on children and armed conflict. For instance, Resolution 2427\footnote{United Nations Security Council. \emph{Resolution 2427 (2018).} UN Doc. S/RES/2427. July 9, 2018}, which will be explained later throughout the study guide, stands as the most cohesive statement on child soldiers to date. Nevertheless, the ratification and enforcement of such frameworks remain limited in reality, particularly in regions where armed groups operate outside of state authority.

The Monitoring and Reporting Mechanism (MRM), established under Security Council Resolution 1612 (2005) , tracks grave violations against children in armed conflict, and its annual Secretary-General reports name all the parties (state and non-state) response.\footnote{United Nations Security Council. \emph{Resolution 1612 (2005).} UN Doc. S/RES/1612. July 26, 2005.} Parties that appear on this “list of shame” are expected to enter into formal Action plans with the UN, committing to end such grave violations. The MRM to date, has remained one of the most operationally fundamental accountability tools available to the international community although the delisting process has shown mixed results.

The challenge of child soldiers remains particularly relevant in current conflicts. In the Democratic Republic of Congo, the Congo, Somalia, South Sudan, the Central African Republic, Myanmar, and Mali, children continue to be recruited and used by both government forces and non-state armed groups.\footnote{United Nations Secretary-General. \emph{Children and Armed Conflict: Report of the Secretary-General.} UN Doc. A/78/842–S/2024/384. May 2024.\url{https://www.un.org/ga/search/view_doc.asp?symbol=S/2024/384}} Multi-actors – where dozens of armed factions may operate simultaneously within a single country – complicates the DDR implementation significantly. The DDR addresses post-conflict transitions, therefore rendering the framework ill-suited to ongoing fighting environments. The committee must explore therefore, how DDR frameworks can be adapted towards dynamic conflict settings that continue without a formal peace agreement.

DiSEC is well-suited to address this issue as DDR programs sit at the heart of disarmament, security, and humanitarian policy (the key mandate of what this committee hopes to achieve)\footnote{United Nations General Assembly. \emph{Disarmament and International Security Committee (First Committee): Mandate and Functions.}\url{https://www.un.org/en/ga/first/}}. Delegates must take into consideration approaches to strengthen accountability for parties that recruit children, close the gap between enforcement in existing legal frameworks, ensure sustainable and adequate funding mechanisms for transition programs, and develop key frameworks that serve all categories of children affected by armed conflict. The question facing this committee is not simply how to free children from armed groups, but rather, how to guarantee them a safe and secure future, and how the international community can build and deliver on that promise.

\subsection{History of the Topic}

\subsubsection{Origins}

The use of children in armed conflicts is not as recent as it may seem. Children have been involved in war contexts since ancient and medieval times. A good example is that of Spartan boys who were taken away from their families at the age of seven to join a military training program referred to as the “agoge” which effectively prepared them to become soldiers once they were older\footnote{History Extra. \emph{Innocence under Fire: How Children were Used in Warfare of the Past}. 2021. \url{https://www.historyextra.com/period/first-world-war/how-children-used-warfare-fighting-child-soldiers/}}. This practice persisted throughout the Modern Era since there is historical evidence that Napoleon resorted to the recruitment of children throughout the final stages of the Napoleonic Wars in 1814 when faced with invasion by Allied forces\footnote{Graeber Daniel. \emph{Child Soldiers: A Timeless Struggle}. Foreign Policy Association. \url{https://fpa.org/child-soldiers-a-timeless-struggle/}}. Later on, about 250,000 boys aged below 18 enlisted themselves in the army to fight for Britain throughout the First World War. Most of them were driven by nationalist narratives or by a desire to escape poverty at home\footnote{\emph{The Teenage Soldiers of World War One}. BBC News.  2014. \url{https://www.bbc.com/news/magazine-29934965}}. Something similar happened in World War Two, where Nazi Germany used slightly untrained and underaged members of the “Hitler Youth” group to fight the Allies, particularly during the battle of Berlin of 1945\footnote{\emph{Hitler Youth: How the Third Reich Used Children’s Organisations to Wage War}. History Extra. 2020. \url{https://www.historyextra.com/period/second-world-war/hitler-youth-children-history-soldiers-fight-jojo-rabbit-film-ww2/}}.

\subsubsection{Modern Child Soldier Crisis}

However, what is known as the “Modern Child Soldier Crisis” only began throughout the 20th Century. It is during this period that various civil wars mainly taking place mainly in Africa but also in other regions made the practice of child soldier use and recruitment more visible. Throughout Sierra Leone’s Civil War taking place between 1991 and 2002, approximately  50\% of the children recruited by rebel forces were abducted at the age of 15 or below\footnote{Betancourt Theresa S. et al. \emph{Sierra Leone’s Child Soldiers: War Exposures and Mental Health Problems}. Gender National Library of Medicine. \url{https://pmc.ncbi.nlm.nih.gov/articles/PMC3124662/}}. In Liberia, children have played a central role in government and rebel armies since the nation’s civil war dating back to the 1990s\footnote{Pan Esther. \emph{Liberia: Child Soldiers}. Council on Foreign Relations. 2005.  \url{https://www.cfr.org/backgrounders/liberia-child-soldiers}}. Similarly, the International Labour Office estimates that, in the Democratic Republic of Congo, the Alliance of Democratic Forces for the Liberation of the Congo (AFDL) recruited about 10,000 children aged between 7 and 16 between the years 1996 and 1997\footnote{International Labour Office. \emph{Wounded Childhood: The Use of Children in Armed Conflict in Central Africa}. 2003. \url{https://www.ilo.org/sites/default/files/wcmsp5/groups/public/@ed_emp/@emp_ent/@ifp_crisis/documents/publication/wcms_116566.pdf}}. In Latin America, particularly in Mexico, children have also historically been used by drug cartels who promise them status and a sense of belonging in exchange for their participation in hostile environments\footnote{Child Soldiers International. \emph{Child Soldiers Global Report 2001 - Mexico}. 2001 \url{https://www.refworld.org/reference/annualreport/cscoal/2001/65235}}.

\subsubsection{Recent Events}

More recently, Human Rights Watch recalled how the UN had been able to verify 2,138 violations against minors in armed conflict taking place in Myanmar in 2024. Among these violations was the recruitment of children into the military\footnote{Human Rights Watch. \emph{Myanmar: Stop Recruitment, Use of Child Soldiers}. 2025. \url{https://www.hrw.org/news/2025/06/20/myanmar-stop-recruitment-use-of-child-soldiers}}. Likewise, it has been documented that over 620 children were recruited by armed criminal groups in Colombia in 2024\footnote{\emph{Kids on the Front Lines: Stopping Child Recruitment in Colombia}. International Crisis Group. Colombia. 2026. \url{https://www.crisisgroup.org/brf/latin-america-caribbean/colombia/b55-kids-front-lines-stopping-child-recruitment-colombia}}. As it will be documented later throughout this study guide, these practices have remained in place in some of the most recent armed conflicts, such as the Russia-Ukraine or Israel-Palestine Conflicts.

\subsection{Past Actions}

The table below shows the various existing treaties and conventions which have sought to protect children from becoming involved in armed conflict:

\begin{guidetable}[Y{0.32}Y{1.12}Y{0.96}Y{1.6}]{4}
  \textbf{Year} & \textbf{Treaty/Convention} & \textbf{Entity} & \textbf{Relevance} \\
  \midrule
  1977 & Protocol Additional to the Geneva Conventions of 12 August 1949, and relating to the Protection of Victims of International Armed Conflicts (Protocol I). & Diplomatic Conference on the Reaffirmation and Development of International Humanitarian Law Applicable in Armed Conflicts. & Article 77(2) of Protocol I sets forth that children below 15 years old shall not be recruited to partake in hostilities\footnote{Diplomatic Conference on the Reaffirmation and Development of International Humanitarian Law Applicable in Armed Conflicts. \emph{Protocol Additional to the Geneva Conventions of 12 August 1949, and relating to the Protection of Victims of International Armed Conflicts (Protocol I)}. 1977. \url{https://ihl-databases.icrc.org/en/ihl-treaties/api-1977/article-77?activeTab=}}. \\
  1989 & Convention on the Rights of the Child (CRC). & Adopted by the UN General Assembly through Resolution 44/25. & Article 38 establishes that children should not be part of armed conflicts, and that those who have been involved in them should obtain assistance in order to recover and be able to fully reintegrate themselves back into society\footnote{UN General Assembly resolution 44/25. \emph{Convention on the Rights of the Child}. 1989. \url{https://www.ohchr.org/en/instruments-mechanisms/instruments/convention-rights-child}}. \\
  1990 & African Charter on the Rights and Welfare of the Child. & African Union. & Article 22 calls upon member states to refrain from recruiting children to take part in armed conflicts\footnote{African Union. \emph{African Charter on the Rights and Welfare of the Child}. Ethiopia. 1990. \url{https://au.int/sites/default/files/treaties/36804-treaty-african_charter_on_rights_welfare_of_the_child.pdf}}. \\
  1998 & Rome Statute. & International Criminal Court. & Article 8(2)(e)(7) defines the conscription or use of children in armed conflicts as a war crime\footnote{International Criminal Court. \emph{Rome Statute}. 1998. \url{https://www.icc-cpi.int/sites/default/files/2024-05/Rome-Statute-eng.pdf}}. \\
  1999 & Worst Forms of Child Labour Convention (No.182). & International Labour Organization. & Article 3(a) recognises “forced or compulsory recruitment of children for use in armed conflict” as one of the “worst forms of child labour”\footnote{International Labour Organization. \emph{Worst Forms of Child Labour Convention (No.182)}. 1999. \url{https://normlex.ilo.org/dyn/nrmlx_en/f?p=NORMLEXPUB:12100:0::NO::P12100_ILO_CODE:C182}}. \\
  2000 & Optional Protocol to the Convention on the Rights of the Child (OPAC) on the involvement of children in armed conflict. & Adopted by the UN General Assembly under Resolution A/RES/54/263. & States that signed this protocol promised not to send children aged below 18 to fight in armed conflicts, as well as to provide physical and psychological assistance to children who had been involved in hostilities\footnote{General Assembly resolution A/RES/54/263. \emph{Optional Protocol to the Convention on the Rights of the Child on the involvement of children in armed conflict}. 2000. \url{https://www.ohchr.org/en/instruments-mechanisms/instruments/optional-protocol-convention-rights-child-involvement-children}}. \\
\end{guidetable}

In addition, the following table summarises some of the most relevant resolutions on the topic of child soldiers adopted by the United Nations Security Council (UNSC), which could be taken as a reference by this committee:

\begin{guidetable}[Y{0.32}Y{0.6}Y{2.08}]{3}
  \textbf{Year} & \textbf{Resolution} & \textbf{Relevance} \\
  \midrule
  1999 & Resolution 1261 & Denounced the recruitment of children in armed conflicts, acknowledging the harmful long-term consequences that this encompasses\footnote{United Nations Security Council. \emph{Resolution 1261}. 1999. \url{https://www.securitycouncilreport.org/atf/cf/\%7B65BFCF9B-6D27-4E9C-8CD3-CF6E4FF96FF9\%7D/CAC\%20SRES\%201261.pdf}}. \\
  2001 & Resolution 1379 & Primarily asked the Secretary-General to identify the parties to armed conflict that recruited or used child soldiers\footnote{United Nations Security Council. \emph{Resolution 1379}. 2001. \url{file:///C:/Users/Usuario/Downloads/S_RES_1379(2001)-EN.pdf}}. \\
  2005 & Resolution 1612 & Condemned the use and recruitment of child soldiers and endorsed the Secretary-General’s Monitoring and Reporting Mechanisms (MRM) action plan aimed at collecting information on armed groups and governments that committed abuses against children during armed conflicts\footnote{United Nations Security Council. \emph{Resolution 1612}. 2005. \url{file:///C:/Users/Usuario/Downloads/S_RES_1612(2005)-EN\%20(1).pdf}}. \\
  2014 & Resolution 2143 & Acknowledged the importance of demobilization, rehabilitation and reintegration programs for children in armed conflicts\footnote{United Nations Security Council. \emph{Resolution 2143}. 2014.  \url{file:///C:/Users/Usuario/Downloads/S_RES_2143(2014)-EN.pdf}}. \\
  2018 & Resolution 2427 & Emphasised that former child soldiers should be guaranteed access to schooling in order to avoid chances of them being recruited again in the future as a result of poverty, mentioned the importance of providing them with psychosocial support, and urged for the different needs of girls and boys to be taken into account for the preparation of tailored reintegration programs\footnote{United Nations Security Council. \emph{Resolution 2427}. 2018. \url{https://unscr.com/en/resolutions/doc/2427/}}. \\
\end{guidetable}

Now that a range of relevant treaties, conventions and UN resolutions have been outlined, it is also worth mentioning how efforts to combat the recruitment of children by parties to armed conflicts have also been made through the implementation of a variety of frameworks, guidelines and initiatives.

In 2014, the UN Special Representative of the Secretary-General for Children and Armed Conflict worked with UNICEF to tackle the issue at hand by launching the “Children, Not Soldiers” campaign, which sought to bring about international consensus on the fact that children should not become involved in war. This initiative released and reintegrated thousands of child soldiers\footnote{UN Office of the Special Representative of the Secretary-General for Children and Armed Conflict. \emph{Child Recruitment and Use}. \url{https://childrenandarmedconflict.un.org/en/six-grave-violations/child-soldiers}}. More generally, it is estimated that UNICEF has been able to deliver protection and reintegration assistance to approximately 12,500 former child soldiers\footnote{Haupt Naomi. \emph{Keeping the Spotlight on Africa’s Child Soldiers}. Institute for Security Studies. 2025. \url{https://issafrica.org/iss-today/keeping-the-spotlight-on-africa-s-child-soldiers} that}. Similarly, the United Nations has partnered with the African Union to foster Disarmament, Demobilization and Reintegration (DDR) Programs focused on removing weapons from child soldiers, releasing them from working for armed groups, and reintegrating them back into society through education and psychological support\footnote{United Nations. \emph{Disarmament, Demobilization, Reintegration (DDR) \& Weapons and Ammunitions Management (WAM) with human rights approach}. United Nations Office to the African Union. \url{https://unoau.unmissions.org/en/disarmament-demobilization-reintegration-ddr-weapons-and-ammunitions-management-wam-with}}.

In 2007, the UN also launched the Principles and Guidelines on Children Associated with Armed Forces or Armed Groups, informally referred to as the Paris Principles, which are a set of guidelines created to prevent child soldiers from being recruited, guarantee their release once recruited, and reintegrate them back into society after they have been released\footnote{United Nations. \emph{Principles and Guidelines on Children Associated with Armed Forces or Armed Groups}. 2007. \url{https://www.unicef.org/mali/media/1561/file/parisprinciples.pdf}}. The Vancouver Principles on Peacekeeping and the Prevention of the Recruitment and Use of Child Soldiers of 2017, which are a set of political commitments aimed at eliminating the use of child soldiers during UN peacekeeping operations, are also relevant since they advocate that peacekeepers must be trained on how to protect minors and that violations against children’s rights during peacekeeping operations must be adequately monitored and reported\footnote{Government of Canada. \emph{The Vancouver Principles on Peacekeeping and the Prevention of the Recruitment and Use of Child Soldiers}. 2017. \url{https://www.canada.ca/en/department-national-defence/corporate/reports-publications/vancouver-principles/introduction.html}}.

Likewise, while future sections in this study guide will provide a more detailed overview of the different regional efforts aimed at addressing the problem of child soldiers, it is worth noting how the European Union’s adoption of the Guidelines on Children and Armed Conflict in 2003 stands out as a key initiative which supports the prevention of child soldier use and recruitment through DDR programs\footnote{European Union. \emph{EU Guidelines on Children and Armed Conflict}. 2003. \url{https://www.eeas.europa.eu/sites/default/files/documents/2024/EEAS-EU-Guidelines-CAAC.pdf}}. Likewise, the Kathmandu Declaration on the Use of Children as Soldiers is also a relevant statement which called for the implementation of measures such as the denial of arms to groups using child soldiers. This was signed by 30 different countries at the Asia-Pacific Conference held in Nepal in the year 2000\footnote{Human Rights Watch. \emph{Kathmandu Declaration on the Use of Children as Soldiers}. \url{https://www.hrw.org/news/2000/05/18/kathmandu-declaration-use-children-soldiers}}.

Finally, the Coalition to Stop the Use of Child Soldiers used to stand out as an active global coalition composed of human rights and humanitarian organizations that sought to monitor, denounce, and prevent the use of children in war contexts. Its main activities included the drafting of reports that document the recruitment of child soldiers across the world and the development of campaigns seeking to raise awareness about the problem as well as to advocate for the implementation of stronger and more effective international laws aimed at deterring armed groups and governments from using child soldiers\footnote{Amnesty International. \emph{Coalition to Stop the Use of Child Soldiers}. 2004. \url{https://www.amnesty.org/es/wp-content/uploads/2021/08/act760022004en.pdf}}.
